\setcounter{page}{1}
\setcounter{figure}{0}
\setcounter{chapter}{4}


\chapter{Movimiento parabólico}
\label{cap:cap5}
\vspace{-15ex}

\begin{flushright}
\linea{2.0}
\end{flushright}
%\fancyfoot{}
%\fancyfoot[LE,RO]{}
\raggedright
\begin{enumerate}
%\item Una intr\'epida nadadora se lanza desde un risco con un impulso horizontal $\VEC{v}_0=v_0 [\HAT{\iota}]$. ?`Qu\'e rapidez m\'inima debe tener al saltar de lo alto del risco para no chocar con la saliente en la base, que tiene una anchura de 1.75 m y est\'a a 9.00 m por debajo del borde del risco?
%\item Se dispara un proyectil desde el nivel del suelo con una velocidad inicial de 50.0 m/s a 60.0$^{\circ}$ por encima de la horizontal sin que sufra resistencia del aire. $a)$  Determine las componentes horizontal y vertical de la velocidad inicial del proyectil. $b)$ ?`Cu\'anto tarda el proyectil en alcanzar su punto m\'as alto? $c)$ Calcule la altura m\'axima que alcanza respecto al suelo. $d)$ ?`Qu\'e tan lejos del punto de lanzamiento cae el proyectil al suelo? $e)$ Determine las componentes horizontal y vertical de la aceleraci\'on y velocidad del proyectil en su punto de altura m\'axima y cuando impacta el suelo.

\begin{comment}
\item \label{item:martillo} Un trabajador se encuentra en un pozo de 2,80 m de profundidad y a 3,50 m de la pared del pozo. Su jefe, que est\'a fuera del pozo a una distancia L del borde, le pide que le lance el martillo. El martillo sale de la mano del trabajador a 1,20 m del
fondo del pozo con un \'angulo de $38,0^{0}$ (ver Figura \ref{fig:martillo}).
\begin{figure}[h!]
\centering
\includegraphics[width=5in]{../figuras/martillo.png}
\caption{Lanzamiento de martillo desde pozo. Ejercicio \ref{item:martillo}.}
\label{fig:martillo}
\end{figure}

\begin{enumerate}
\item ?`Cu\'al es la rapidez inicial m\'inima necesaria para que el martillo pueda salir del pozo (rozando el borde)?
\item ?`Qu\'e distancia $L$ debe haber entre el borde del pozo y el jefe para que el martillo no lo golpee antes de que caiga al suelo?
\item ?`Cu\'anto tiempo tarda el martillo para tocar el suelo luego de que sale del pozo?
\end{enumerate}  
\end{comment}

\item Una persona se encuentra en el borde de un \href{https://viajes.nationalgeographic.com.es/a/8-acantilados-costeros-europa-que-debes-ver-vez-vida-al-menos_9637}{\textcolor{c2}{acantilado}} (paredón que cae verticalmente) sobre el mar, a una altura de 300 ft. Lanza una piedra con un \'angulo sobre la horizontal de 30$^\circ$. La piedra tarda 5,00 s en llegar al mar. Calcule
\begin{enumerate}
\item la rapidez con que inicialmente fue lanzada la piedra,
\item la distancia horizontal que recorre la piedra hasta que cae al mar,
\item la altura m\'axima sobre el nivel del mar que alcanza la piedra.
\end{enumerate}


\item  Un helicoptero viaja a una altura constante de 8 mil pies de altura, a una velocidad de 450 km/h cuando libera un paquete. Determine
\begin{enumerate}
\item el tiempo que tarda el paquete en impactar el suelo, 
\item la distancia horizontal que recorre el paquete.
\item la velocidad (magnitud y direcci\'on) a la que impactar\'a el paquete el suelo. 
\end{enumerate}
\newpage

\item \textsc{\textcolor{c2}{[2P/2S/2014]}}  \label{item:amigos} Tres amigos temerarios corren a trav\'es de los tejados de la ciudad. Los tres est\'an corriendo con una rapidez de 5 m/s cuando llegan a un espacio vac\'io entre dos edificios. El espacio vac\'io tiene 4 m de anchura y presenta un desnivel de 3 m por debajo del tejado del primer edificio (ver Figura \ref{fig:2p2s2014}). El primero salta con un \'angulo de 45$^{\circ}$ sobre la horizontal y cruza el hueco con facilidad. El segundo amigo salta horizontalmente porque piensa que es lo mejor. En estas circunstancias determine 

\begin{figure}[h!]
\begin{center}
\includegraphics[width=3.5in]{../figuras/amigos_temerarios.png}
\caption{Amigos temerarios. Ejercicio \ref{item:amigos}}
\label{fig:2p2s2014}
\end{center}
\end{figure}
\begin{enumerate}
\item ?`A qu\'e distancia del borde del segundo edificio llega el primer amigo?% \\
%\textcolor{blue}{El salto del primer amigo est\'a caracterizado por $$v_0=5\units{m/s},\quad \mbox{y} \qquad \theta=45^{\circ}.$$ Podemos resolver la ecuaci\'on de la trayectoria para encontrar en que valor de ``$x$" se cumple $$y_1(x)=-3\units{m},$$ es decir, $$y_{1}(x)=x\cdot \tan (45^{\circ})-\frac{g}{2\cdot(5\units{m/s})^2 \cos ^2 (45^{\circ})}\cdot x^2 =-3\units{m};$$ lo que resulta en $x=4.3\units{m}.$ \resp{}{El primer amigo cae a 30 cm del borde del segundo edificio.}}
\item  ?`Lograr\'a el segundo amigo cruzar el espacio vac\'io entre los edificios? %\\ \textcolor{blue}{El salto del segundo amigo est\'a caracterizado por $$v_0=5\units{m/s},\quad \mbox{y} \quad \theta=0^{\circ}.$$ Podemos resolver la ecuaci\'on de la trayectoria para encontrar para que valor de ``$x$" se cumple $$y_2(x)=-3\units{m}, $$ lo que resulta en en $x=3.9\units{m}$. \resp{.9\textwidth}{El segundo amigo no cruza el espacio entre los edificios, pues le faltaron 10 cm.}}
\item  ?`Qu\'e condici\'on debe cumplir el salto del tercer amigo para que logre caer en el segundo edificio?  %\\ \textcolor{blue}{El tercer amigo debe saltar con un \'angulo $\alpha$ que satisfaga que: \\ \resp{}{$-3\units{m}=\tan\alpha \cdot (4\units{m})^2 -\displaystyle \frac{g}{2\cdot(5\units{m/s})^2 \cos^2 \alpha}\cdot (4\units{m})^2 $}}
\end{enumerate}

\newpage

\item \label{item:volcan} En una violenta erupci\'on, el volc\'an Turrialba lanza una roca con una rapidez de 200 km/h y una elevaci\'on de $\theta_0=23^\circ$. El cr\'ater del volc\'an se encuentra a 3340 m de altura (ver Figura \ref{fig:volcan}).

\begin{figure}[h!]
    \centering
\includegraphics[width=4in]{../figuras/volcan.png}
    \caption{Roca expulsada en erupci\'on volc\'anica. Ejercicio \ref{item:volcan}}
    \label{fig:volcan}
\end{figure}

Determine

\begin{enumerate}
\item  la distancia $d$ a la que cae la roca de la base del volc\'an,
\item la altura m\'axima (respecto al suelo) que alcanza la roca,
\item ?`cu\'anto tarda la roca en impactar el suelo?
\item la velocidad (magnitud y direcci\'on) de la roca al impactar el suelo.
\end{enumerate}

%\item Los meteoros son rocas de hasta unos cuantos metros que ingresan a la atm\'osfera terrestre, y debido a la fricci\'on que experimentan se desintegran poco a poco. Si las condiciones de ingreso a la atm\'osfera son adecuadas, logran impactar la superficie terrestre sin desintegrarse por completo, en cuyo caso se les da el nombre de \textit{meteorito}. Un meteorito que ingres\'o a la atm\'osfera con un \'angulo de $\theta=30^{\circ}$, se descubri\'o a 150 km horizontalmente de su punto de ingreso. $a)$ Haga un diagrama que esquematice la trayectoria del meteorito una vez que ingresa a la atm\'osfera. Justifique. $b)$ Determine la rapidez inicial con que ingres\'o el meteorito a la atm\'osfera. %$c)$ Determine la velocidad (magnitud y direcci\'on) con la que el meteorito impact\'o el suelo. \\

%\begin{center}\fbox{\begin{minipage}{.8\textwidth} Considere que la atm\'osfera terrestre tiene una altura de 100 km, y dado que la roca alcanz\'o impactar el suelo, desprecie la fricci\'on de la atm\'osfera.\end{minipage}}
%\end{center}


%\item Mantenemos el equilibrio, al menos en parte, gracias a la endolinfa del o\'ido interno. El giro de la cabeza desplaza este l\'iquido, produciendo mareos. Suponga que un bailar\'in est\'a girando muy r\'apido a 3.0 revoluciones por segundo alrededor de un eje vertical que pasa por el centro de su cabeza. Aun cuando la distancia var\'ia de una persona a otra, el o\'ido interno se encuentra aproximadamente a 7.0 cm del eje de giro. ?`Cu\'al es la aceleraci\'on radial (en m/s$^{2}$) de la endolinfa?
%\item Un modelo de rotor de helic\'optero tiene cuatro aspas, cada una de 3.40 m de longitud desde el eje central hasta la punta. El modelo gira en un t\'unel de viento a 550 rpm. $a)$ Qu\'e rapidez lineal tiene la punta del aspa en m/s? $b)$?`Qu\'e aceleraci\'on radial tiene la punta del aspa?
%\item Un sat\'elite de telecomunicaciones gira alrededor de la Tierra en una \'orbita circular de 35800 km de radio. Si el per\'iodo orbital del sat\'elite es el mismo que el per\'iodo rotacional de la Tierra sobre su propio eje, calcule $a)$ la velocidad tangencial del sat\'elite, $b)$ la aceleraci\'on radial del sat\'elite.
%\item Un jet vuela en picada para realizar un maniobra. Al llegar a 500 m del suelo cambia su direcci\'on, haciendo una curva con un radio de 350 m. De acuerdo a pruebas m\'edicas, los pilotos pierden la conciencia a una aceleraci\'on de 5.5$g$ ?`A qu\'e rapidez de descenso perder\'ia la conciencia el piloto al realizar esta maniobra?
%\item Una banda m\'ovil en un aeropuerto se mueve a 1.0 m/s y tiene 35.0 m de largo. Si una mujer entra en un extremo y camina a 1.5 m/s en relaci\'on con la banda m\'ovil, ?`cu\'anto tiempo tardar\'a en llegar al otro extremo si camina $a)$ en la misma direcci\'on que se mueve la banda? $b)$ ?`en la direcci\'on opuesta?
%\item Un r\'io fluye al sur con rapidez de 2.0 m/s. Un hombre cruza el r\'io en una lancha de motor con velocidad relativa al agua de 4.2 m/s al este. El r\'io mide 800 m de ancho. $a)$ ?`Qu\'e velocidad (magnitud y direcci\'on) tiene la lancha relativa a la Tierra? $b)$ ?`Cu\'anto tiempo tarda en cruzar el  r\'io? $c)$ ?`A qu\'e distancia al sur de su punto de partida llegar\'a a la otra orilla?
%\item Los gansos canadienses viajan principalmente en direcci\'on norte-sur por mucho m\'as de mil kil\'ometros en algunos casos, viajando a velocidades hasta de 100 km/h aproximadamente. Si una de estas aves vuela a 100 km/h en relaci\'on al aire, pero hay un viento de 40 km/h que sopla de oeste a este, $a)$?`a qu\'e \'angulo en relaci\'on con la direcci\'on norte-sur deber\'ia volar esta ave  de modo que viaje directamente hacia el sur en relaci\'on al suelo? $b)$ ?`Cu\'anto tiempo le tomar\'a al ganso cubrir una distancia de 500 km de norte a sur?
\end{enumerate}
